
\documentclass{article} %210 mm × 297 mm
\usepackage[utf8]{inputenc}
\usepackage[a4paper,left=1.5cm, right=1.5cm, top=2cm, bottom=2cm]{geometry}
\usepackage{graphicx}
\usepackage{systeme}          % for Solution 3
\usepackage{array}  
%\usepackage{blindtext}
\usepackage{polynom}
\usepackage{multirow}
\usepackage{mhchem}
\usepackage[normalem]{ulem}
\usepackage{url}
\usepackage{subfigure}
\usepackage{amsfonts,latexsym} % para tener disponibilidad de diversos simbolos
\usepackage{enumerate}
%\usepackage{booktabs}
\usepackage{float}
%\usepackage{threeparttable}
\usepackage{array,colortbl}
%\usepackage{ifpdf}
%\usepackage{rotating}
\usepackage{stfloats}
\usepackage{url}
\usepackage{listings}
\usepackage{fancyhdr}
\usepackage{biblatex}
\usepackage{color}
\addbibresource{sources.bib}
%\usepackage{hyperref}
%\usepackage{wrapfig}
\usepackage{array}
%\usepackage{multirow}
%\usepackage{adjustbox}
%\usepackage{nccmath}
\usepackage[dvipsnames]{xcolor}
\usepackage{tcolorbox}
\newtcolorbox{mybox}{colback=white,
colframe=black}
\newcommand{\titulo}{Abgabe 10; Diskrete Mathematik}
\newcommand{\nombre}{Liliana Sanfilippo}
\newcommand{\FCE}{\textbf{WiSe 2024/2025}}

 

\usepackage{fancyhdr}
\pagestyle{fancy}
\fancyhead{}
\fancyfoot{}
\fancyhead[LO]{\small\nombre}
\fancyhead[RE]{\small\titulo}
\fancyhead[LO]{\small\FCE}
\fancyfoot[RO,LE] {\thepage}
\usepackage[onehalfspacing]{setspace}

\setcounter{page}{1} %Número de la página, la editorial lo cambiará
\begin{document}
\begin{tabular}{p{12cm}}
{{\Large\bf\titulo}}\\
\vspace{0.2cm}\\
 {\large\nombre}\\
 \vspace{0.1cm}
\end{tabular}

\section*{Präsenzaufgaben}
\begin{enumerate}
    \item Dividieren Sie mit Rest
    \[
    x^5 + x^4 + 2x^3 + x^2 + 4x + 2
    \]
    durch
    \[
    x^2 + 2x + 3
    \]
    im Polynomring \( \mathbb{Z}/5\mathbb{Z}[x] \).
    \begin{mybox}
    \[
    r(x) = x^5 + x^4 + 2x^3 + x^2 + 4x + 2
    \]
    \[
    g(x) = x^2 + 2x + 3
    \]
    Leitkoeffizient von $r(x): 1$ \\
    Leitkoeffizient von $g(x): 1$ \\
    \[
    \frac{1}{1} \mod 5 = 1 
    \]
     \[
    t(x)_1 = x^3
    \]
    \[
    q(x) = x^3
    \]
     \begin{align*}
        r(x) & = (x^5 + x^4 + 2x^3 + x^2 + 4x + 2) - (x^3 \cdot (x^2 + 2x + 3)) \mod 5 \\
        r(x) & = (x^5 + x^4 + 2x^3 + x^2 + 4x + 2) - (x^5 + 2x^4 + 3x^3) \mod 5 \\
        r(x) & = (-x^4 - 1x^3 + x^2 + 4x + 2) \mod 5 \\
        r(x) & = 4x^4 +4x^3 + x^2 + 4x + 2
    \end{align*}
    \hrule 
    \[
    \frac{-1}{1} \mod 5 = -1 \mod 5 = 4
    \]
    \[
    t(x)_2 = 4x^2
    \]
    \[
    q(x) = x^3 + 4x^2
    \]
     Leitkoeffizient von $r(x): -1$ \\
    Leitkoeffizient von $g(x): 1$ \\
    \begin{align*}
        r(x) & = (4x^4 +4x^3 + x^2 + 4x + 2) - (4x^2 \cdot (x^2 + 2x + 3)) \mod 5 \\
        r(x) & = (4x^4 +4x^3 + x^2 + 4x + 2) - (4x^4 + 8x^3 + 12x^2) \mod 5\\
        r(x) & = (4x^4 +4x^3 + x^2 + 4x + 2) - (4x^4 + 3x^3 + 2x^2) \\
        r(x) & = (x^3 - x^2 + 4x + 2)
    \end{align*}
    \hrule 
    \end{mybox}
    
    
    \begin{mybox}
    Leitkoeffizient von $r(x): 1$ \\
    Leitkoeffizient von $g(x): 1$ \\
    \[
    \frac{1}{1} \mod 5 = 1 \mod 5 = 1
    \]
    Gradunterschied: 1
    \[
    t(x)_3 = x
    \]
    \[
    q(x) = x^3 + 4x^2 + x
    \]
    \begin{align*}
        r(x) & = (x^3 - x^2 + 4x + 2) - (x \cdot (x^2 + 2x + 3)) \mod 5 \\
        r(x) & = (x^3 - x^2 + 4x + 2) - (x^3 + 2x^2 + 3x) \mod 5 \\
        r(x) & = ( - 3x^2 + x + 2) \mod 5 \\
        r(x) & = 2x^2 + x + 2
    \end{align*}
   \sout{ Lösung: 
    $q(x) = x^3 + 4x^2 + x$ mit  
    $r(x)  = ( - 3x^2 + x + 2) $} Da muss man noch einen Schrott weiter, damit der Rest einen Grad kleiner hat als $g(x)$
    \hrule 
    Leitkoeffizient von $r(x): 2$ \\
    Leitkoeffizient von $g(x): 1$ \\
    \[
    \frac{2}{1} \mod 5 = 2 \mod 5 = 2
    \]
    Gradunterschied: -1
    \[
    t(x)_4 = 2
    \]
    \[
    q(x) = x^3 + 4x^2 + x + 2
    \]
    \begin{align*}
        r(x) & = (2x^2 + x + 2) - (2 \cdot (x^2 + 2x + 3)) \mod 5 \\
       r(x) & = (2x^2 + x + 2) - (2x^2 + 4x + 6) \mod 5 \\
       r(x) & = (-3x -4) \mod 5 \\
       r(x) & = 2x + 1
    \end{align*}
    Lösung: 
    $q(x) = x^3 + 4x^2 + x + 2$ mit  
    $r(x)  = 2x + 1 $
    \end{mybox}
        \item Finden Sie alle monischen (normierten) irreduziblen Polynome in \( \mathbb{Z}/3\mathbb{Z}[x] \) vom Grad höchstens 2.
        \begin{mybox}
        Wir suchen Polynome der Form $x^2+ax+b$ (Grad 2), $x+b$ (Grad 1) und $b$ (Grad 0) mit $a,b \in \mathbb{Z}/3\mathbb{Z}$. \\
        Da $\mathbb{Z}/3\mathbb{Z} = {0,1,2}$ beinhaltet es folgende Polynome mit diesen Eigenschaften: 
       \[
       \begin{array}{cccc}
        x^2 &  & & \text{reduzibel}\\
        x^2 &  & + 1 \\ 
        x^2 &  & + 2 \\
        x^2 & + x & \\ 
        x^2 & + x & + 1 \\ 
        x^2 & + x & + 2 \\ 
        x^2 & + 2x &  \\ 
        x^2 & + 2x & + 1 \\ 
        x^2 & + 2x & + 2 \\
        & x \\
        &  x & + 1 \\
        & x & +2 \\
        & & 1 & irreduzibel\\
        & & 2
       \end{array}
       \]
       Ich verstehe nicht, wie ich auf eine (einigermaßen zeit-freundliche Art) berechnen kann, wie man berechnet, ob ein Polynom reduzibel ist. 
        \end{mybox}
\end{enumerate}


\section*{Aufgaben zur Abgabe}

\subsection*{Aufgabe 1 (4 Punkte)}
Finden Sie den größten gemeinsamen Teiler der Polynome
\[
x^6 + x^5 + x^4 + x^2 + x + 1 \quad \text{und} \quad x^9 + x^8 + x + 1
\]
im Polynomring \( \mathbb{Z}/2\mathbb{Z}[x] \).
\begin{mybox}
%Um den ggT zu berechnen, muss ich erst einmal Polynomdivision machen. Ich habe erst das Polynom mit dem höheren Grad durch das mit dem niedrigeneren Grad geteilt - muss man das so machen?
%\[
%(x^6 + x^5 + x^4 + x^2 + x + 1 ) = (x^9 + x^8 + x + 1) \cdot ( x^{-3} + x^{-5}
%\]
%\[
%\begin{array}{ccccccccccccc}
%    & x^6   & +x^5  &+x^4   &+x^3   &+x^2   &+x &+1 &  \\
%-(  & x^6   & +x^5  &       &       &       &   &   & +x^{-2}&+x^{-1} & ) && | \text{ Rest multipliziert mit }  x^{-3}  \\ \hline
%    &       &       &x^4    &+x^3   &+x^2   &+x &+1 & -x^{-2}&-x^{-1} \\
%-(  &       &       &x^4    &+x^3   &       & & & &        &+x^{-4}& -x^{-5}) & | \text{ Rest multipliziert mit }  x^{-5} \\ \hline
%\end{array}
%\]
%\begin{align*}
%    (x^9 + x^8 + x + 1)&  = (x^6 + x^5 + x^4 + x^2 + x + 1 ) \cdot (x^3 - x +1) -x^5 -2x^4 +x \\
%    (x^9 + x^8 + x + 1) +x^5 +x^4 -x &  = (x^6 + x^5 + x^4 + x^2 + x + 1 ) \cdot (x^3 - x +1) \\
%    x^9 + x^8 +x^5 +x^4 + 1  &  = (x^6 + x^5 + x^4 + x^2 + x + 1 ) \cdot (x^3 - x +1) \\
%    x^9 + x^8 +x^5 +x^4 + 1  &  = (x^6 + x^5 + x^4 + x^2 + x + 1 ) \cdot (x^3 - x +1)
%\end{align*}
Das mit den Vorzeichen in $\mathbb{Z}/2\mathbb{Z}[x]$ finde ich sehr verwirrend. Habe ich das bei dieser Polynomdivision richtig gemacht? 
\[
\begin{array}{cccccccccccc}
    & x^9   & + x^8 & & & & && &   + x & + 1 &  \\
-(  & x^9   & + x^8 & + x^7 &   & +  x^5 & + x^4 & + x^3& ) &   & & | \text{ Rest multipliziert mit } x^3 \\ \hline
    &       &       &  x^7 &       & + x^5 & + x^4 & +x^3  &       &+x & +1 &  \\
-(  &       &       &  x^7 &+x^6   &  +x^5 &       & +x^3  & +x^2  &+x & )& | \text{ Rest multipliziert mit } x \\ \hline
    &       &       &       & x^6   &       & +x^4  &       & +x^2  & +x& +1 \\
-(  &       &       &       & x^6   &  +x^5 & + x^4 &       & +x^2  & +x& +1) & | \text{ Rest multipliziert mit } 1 \\ \hline
&           &       &       &       & x^5  & +x^4 &       &       
\end{array}
\]

\begin{align*}
    r(x) & = x^9 + x^8 + x + 1 \\
    g(x) &= x^6 + x^5 + x^4 + x^2 + x + 1 
\end{align*}
\end{mybox}
\begin{mybox}
 Schritt 1. im euklidischen Algorithmus: 
\[ 
(x^9 + x^8 + x + 1) = (x^6 + x^5 + x^4 + x^2 + x + 1 ) \cdot (x^3 + x +1) + (x^5 + x^4)
\]
Als nächstes berechnen: 
\[
(x^6 + x^5 + x^4 + x^2 + x + 1 ) = (x^5 + x^4) \cdot ? + ? 
\]
Also $(x^5 + x^4) / (x^3 + x +1)$ berechnen. 
\[
\begin{array}{ccccccccc}
    & x^6   & +x^5  &+x^4   &+x^3   &+x^2   &+x &+1 &  \\
-(  &x^6    & +x^5 & )  \\\hline
    &       &       &+x^4   &+x^3   &+x^2   &+x &+1 &  \\
\end{array}
\]
 Schritt 2. im euklidischen Algorithmus: 
\[
(x^6 + x^5 + x^4 + x^2 + x + 1 ) = (x^5 + x^4) \cdot x + (x^4 + x^2 + x + 1)
\]
Als nächstes berechnen:
\[
 (x^5 + x^4) = (x^4 + x^2 + x + 1) \cdot ? + ? 
\]

\[
\begin{array}{ccccccccc}
    &x^5&+x^4 \\
 -( &x^5   &   &+x^3   &+x^2 &+x & ) \\ \hline
    & & x^4 & +x^3 & +x^2 & +x
\end{array}
\]
Schritt 3. im euklidischen Algorithmus: 
\[
 (x^5 + x^4) = (x^4 + x^2 + x + 1) \cdot (x) + (x^4+x^3+x^2+x)
\]
Als nächstes berechnen:
\[
(x^4 + x^2 + x + 1) = (x^4+x^3+x^2+x) \cdot ? +? 
\]

\[
\begin{array}{ccccccccc}
  & x^4 & & + x^2 & + x & + 1  \\
 -( & x^4 & +x^3 & +x^2 & +x & ) \\ \hline
 & & x^3  & & & +1
\end{array}
\]
Schritt 4. im euklidischen Algorithmus:
\[
(x^4 + x^2 + x + 1) = (x^4+x^3+x^2+x) \cdot 1 + (x^3 +1) 
\]
Als nächstes berechnen:
\[
(x^4+x^3+x^2+x) = (x^3 +1) \cdot ? + ? 
\]

\[
\begin{array}{ccccccccc}
  & x^4 & + x^3 & + x^2 & + x \\ 
-(& x^4 & & & +x & ) \\ \hline
& & x^3 + & x^2 
\end{array}
\]
Schritt 5. im euklidischen Algorithmus:
\[
(x^4+x^3+x^2+x) = (x^3 +1) \cdot (x) +  (x^3 +  x^2) 
\]
Als nächstes berechnen:
\[
(x^3 +1) = (x^3 +  x^2)  \cdot ? + ? 
\]
\end{mybox}

\begin{mybox}
\[
\begin{array}{ccccccccc}
  x^3 & & +1 \\
  x^3 & +  x^2 \\ \hline
  x^2 & & +1
\end{array}
\]
Schritt 5. im euklidischen Algorithmus:
\[
(x^3 +1) = (x^3 +  x^2)  \cdot 1 + (x^2 +1) 
\]

Das hat nicht so funktioniert, wie es sollte. Aber ich sitze schon mehr als 2 Stunden an dieser Aufgabe und es wird nicht besser. Daher bleibt es jetzt so. 
\end{mybox}

\subsection*{Aufgabe 2 (4 Punkte)}
Ein Polynom \( f(x) \in k[x] \) heißt monisch, oder normiert, falls der Leitkoeffizient 1 ist. Mit
anderen Worten, \( \deg(f) = n \) und \( a_n = 1 \).

\begin{enumerate}
    \item[(a)] Zeigen Sie, dass es genau \( p \) irreduzible monische Polynome vom Grad 1 in \( \mathbb{Z}/p\mathbb{Z}[x] \) gibt.
    \begin{mybox}
     Ein normiertes Polynom vom Grad 1 in  \( \mathbb{Z}/p\mathbb{Z}[x] \) hat die Form $f(x)=x+a$ für $a \in \mathbb{Z}/p\mathbb{Z}$. Weil $|\mathbb{Z}/p\mathbb{Z}| = p$ und $a$ somit $p$ verschieden viele Werte annehmen kann, gibt es $p$ monische Polynome vom Grad 1 in \( \mathbb{Z}/p\mathbb{Z}[x] \). \\
     Ein Polynom von Grad 1 ist wegen seiner Form $f(x)=x+a$ immer irreduzibel, weil man weder das $x$ noch eine Zahl $a$ weiter zerlegen kann. 
    \end{mybox}
    \item[(b)] Zeigen Sie weiterhin, dass es \( \frac{1}{2}(p^2 - p) \) irreduzible monische Polynome vom Grad 2 in \( \mathbb{Z}/p\mathbb{Z}[x] \) gibt.
    \begin{mybox}
    Ein normiertes Polynom vom Grad 2 in \( \mathbb{Z}/p\mathbb{Z}[x] \) hat die Form $f(x)=x^2+ax+b$ für $a,b \in \mathbb{Z}/p\mathbb{Z}$. \\
    Es gibt jeweils für $a$ und $b$ genau $p$ mögliche Werte und somit $p^2$ viele monische Polynome 2ten Grades in \( \mathbb{Z}/p\mathbb{Z}[x] \). \\
    Von dieser Zahl muss man die reduziblen monischen Polynome abziehen. Da diese sich in zwei irreduzible monische Polynome vom Grad 1 zerlegen lassen würden, gibt es $p$ Möglichkeiten. Man landet bei der Zahl $p^2-p$. \\
    Weil es egal ist, in welcher Reihenfolge die Polynome ersten Grades nach der Zerlegung stehen (siehe $f(x)=(x-r_1)(x-r_2)$ ), müssen wir beachten die Paare $(r_1, r_2)$ nicht doppelt zu zählen. Deshalb teilen wir die Anzahl noch durch 2. 
    \[ \frac{p^2-p}{2} = \frac{1}{2}(p^2 - p)\]
    \end{mybox} 
    \item[(c)] Wie viele irreduzible Polynome vom Grad 3 in \( \mathbb{Z}/p\mathbb{Z}[x] \) gibt es?
    \begin{mybox}
    Monische Polynome vom Grad 3: $p^3$ \\
    Reduzible monische Polynome vom Grad 3: anscheinend $p$
     \[
     \frac{1}{3}(p^3-p)
     \]
    \end{mybox}
\end{enumerate}

Hinweis zu (b): Wie viele monische Polynome vom Grad 2 gibt es? Falls ein solches
Polynom reduzierbar ist, dann faktorisiert es sich als zwei irreduzible monische Polynome
vom Grad 1. Wie viele Möglichkeiten gibt es?

\subsection*{Aufgabe 3 (4 Punkte)}
Sei \( C \) ein linearer Code der Länge \( n \) und Dimension \( k \) mit Minimalabstand \( \delta \). Außerdem
sei \( H \) eine Kontrollmatrix von \( C \).
\begin{enumerate}
    \item[(a)] Zeigen Sie, dass der Minimalabstand \( \delta \) die größte natürliche Zahl \( d \) ist, für die jede
    Menge von \( d - 1 \) Spalten der Kontrollmatrix \( H \) linear unabhängig ist.
    \begin{mybox}
      Beweis durch Widerspruch. \\
      Sei $\delta > d$, dann sind mindestens $\delta$ Spalten von $H$ linear abhängig (weil $\delta > d$ und $d-1$ ist die Grenze, also muss eine Anzahl größer als $\delta -1$ linear abhängig sein). \\
      Dadurch existiert ein Wort mit Hamming-Gewicht kleiner als $\delta$. \\
      \textbf{Widerspruch}, da das Minimalgewicht gleich dem minimalen Hamming-Gewicht ist. 
    \end{mybox}
    \item[(b)] Beweisen Sie die sogenannte Singleton-Schranke
    \[
    \delta \leq n - k + 1
    \]
    \begin{mybox}
    Sei $q$ die Größe des Alphabets\\ 
     Es gibt $q^k$ Codewörter der Länge $n$. Wenn man von jedem Codewort die ersten $\delta-1$ Stellen löscht, müssen die Codewörter weiterhin verschieden sein, weil sie einen Hamming-Abstand von mindestens $\delta$ hatten und haben nun einen von mindestens 1. Außerdem gäbe es weiterhin $q^k$ viele dieser neuen Codewörter. Also $q^{n-\delta +1} = q^k$. Bestehen weitere Einschränkungen für den ursprünglichen Code (mit $q^k$ Wörtern), kann es sein, dass aus dem neuen Code (mit $q^{n-\delta +1}$ Wörtern) nicht alle Wörter auch Wörter im Sinne des alten Codes wären. Daher beschränkt man sich darauf es als Ungleichung zu formulieren:
     \[
     q^{n-\delta +1} \geq q^k   \Rightarrow  n- \delta +1 \geq k  \Leftrightarrow \delta \leq n-k+1
     \]
    \end{mybox}
\end{enumerate}

\subsection*{Aufgabe 4 (4 Punkte)}
Ein linearer Code, für den Gleichheit in der Singleton-Schranke gilt, wird als separabler
Maximal-Distanz-Code (MDS-Code) bezeichnet.
Es sei \( C \) die Menge aller Codewörter der Länge \( n \) mit geradem Gewicht. Zeigen Sie, dass \( C \)
ein linearer MDS-Code ist.
\begin{mybox}
$C$ beinhaltet den Nullvektor, da 0 ein gerades Gewicht ist. Außerdem ist $C$ abgeschlossen unter Addition, da die Summe zweier geraden Gewichte wieder gerade ist, da für die Gewichte von Codewörtern gilt $\partial(a, b) = \partial(a + b, b + b) = \partial(a + b, 0) = w(a + b)$. Also ist $C$ ein linearer Code. \\ 
Das minimalste Gewicht ungleich Null wäre 1, aber das ginge nur, wenn Codewörter auch ungerade wären. Also ist $\delta = 2$. \\ 
Weil der Code nur die Hälfte aller Codewörter der Länge $n$ beinhaltet (die geraden...) gibt es $\frac{2^n}{2} = 2^{n-1}$ viele Wörter in $C$ und die Dimension ist somit $n-1$. 
\begin{align*}
    \delta &  \leq n-k+1\\
     2 &  \leq n- (n-1) +1 \\
     2 &  \leq n- n + 1 +1  \\
     2 &  \leq 2
\end{align*}
Die Singleton-Schranke ist erfüllt und somit ist $C$ ein linearer MDS-Code. 
\end{mybox}

\end{document}